\documentclass[12pt]{article}
\usepackage{amsmath, amssymb}
\usepackage{geometry}
\geometry{margin=1in}
\usepackage{enumitem}

\title{CSE400 – Fundamentals of Probability in Computing}
\author{Lecture 4: Joint Probability and Conditional Probability}
\date{}

\begin{document}
\maketitle

\section*{Covered Portion: Introduction to Probability Theory}

\section*{1. Experiments, Sample Space, and Events}

\subsection*{Definition: Experiment ($E$)}

An experiment is a procedure we perform that produces some result.

\textbf{Example:}

Tossing a coin five times is an experiment denoted as:

\[
E_5
\]

\subsection*{Definition: Outcome ($\xi$)}

An outcome is a possible result of an experiment.

\textbf{Example:}

One outcome of the experiment $E_5$ is:

\[
\xi_1 = HHTHT
\]

\subsection*{Definition: Event}

An event is a certain set of outcomes of an experiment.

\textbf{Example:}

Let $C$ be an event in experiment $E_5$ defined as:

\[
C = \{\text{all outcomes consisting of an even number of heads}\}
\]

\subsection*{Definition: Sample Space ($S$)}

The sample space $S$ is the collection or set of all possible distinct outcomes of an experiment.

\subsection*{Properties}

\begin{itemize}
\item Outcomes are mutually exclusive  
\[
\rightarrow \text{Only one outcome can occur at a time}
\]

\item Outcomes are collectively exhaustive  
\[
\rightarrow \text{No outcome outside the set is possible}
\]
\end{itemize}

\textbf{Example (single coin flip):}

\[
S = \{H,T\}
\]

Cannot get $H$ and $T$ simultaneously (mutually exclusive)

No other result exists (collectively exhaustive)

\subsection*{Continued Properties of Sample Space}

\[
S \text{ is the universal set of outcomes}
\]

Sample space may be:

\begin{itemize}
\item Discrete
\item Countably infinite
\item Continuous
\end{itemize}

\subsection*{Examples of Experiments}

\begin{itemize}
\item Flipping a fair coin once
\item Rolling a cubical die with numbered faces
\item Rolling two dice
\item Flipping a coin until tails occurs
\item Random number generator in interval $[0,1)$
\end{itemize}

\section*{2. Axioms of Probability}

Before assigning probabilities:

Probability is a numerical measure of the likelihood of events, or a function mapping events to numerical values measuring likelihood.

\subsection*{Axiom 1: Boundedness}

For any event $A$:

\[
0 \le Pr(A) \le 1
\]

\subsection*{Axiom 2: Total Probability of Sample Space}

If $S$ is the sample space:

\[
Pr(S) = 1
\]

\subsection*{Axiom 3: Additivity for Mutually Exclusive Events}

If:

\[
A \cap B = \emptyset
\]

then:

\[
Pr(A \cup B) = Pr(A) + Pr(B)
\]

\subsection*{Infinite Version of Axiom 3}

For mutually exclusive events:

\[
A_i \cap A_j = \emptyset \quad \text{for } i \ne j
\]

then:

\[
Pr\left(\bigcup_{i=1}^{\infty} A_i\right) = \sum_{i=1}^{\infty} Pr(A_i)
\]

\section*{3. Corollary from Probability Axioms}

Define:

\[
A_i = \emptyset \quad \text{for all } i > M
\]

\subsection*{Corollary 2.1 (Finite Additivity)}

For a finite number $M$ of mutually exclusive events:

\[
A_i \cap A_j = \emptyset \quad (i \ne j)
\]

then:

\[
Pr\left(\bigcup_{i=1}^{M} A_i\right) = \sum_{i=1}^{M} Pr(A_i)
\]

\subsection*{Notes}

Axiom 3 reduces to Corollary 2.1 when sample space is finite

Axiom 3 is required for infinite sample spaces

\section*{4. Propositions from Probability Axioms}

\subsection*{Proposition 2.1: Complement Rule}

For any event $A$:

\[
Pr(A^c) = 1 - Pr(A)
\]

\subsection*{Proposition 2.2: Monotonicity}

If:

\[
A \subset B
\]

then:

\[
Pr(A) \le Pr(B)
\]

\end{document}

